% ============================================================
%  Chapter 2 — Literature Review and Related Work
% ============================================================
\chapter{Literature Review and Related Work}
\label{chap:lit}

This chapter surveys the body of work most relevant to the problem
of intelligent task scheduling in mobile edge computing environments.
The review is organised thematically, covering edge computing
architectures, classical and heuristic load balancing, machine
learning-based approaches, deep reinforcement learning schedulers,
and predictive techniques that exploit temporal patterns in resource
utilisation. The chapter concludes by identifying the gaps that the
present thesis addresses.

% ─────────────────────────────────────────────────────────
\section{Edge Computing Paradigms}
\label{sec:lit_edge}

The development of edge computing was driven by the inadequacies of
centralised cloud systems in meeting the latency requirements of IoT
applications. By moving compute resources closer to users and data
sources, edge architectures reduce end-to-end latency and conserve
WAN bandwidth. Three primary architectural models have emerged.

\subsection{Cloudlets}

The cloudlet concept was formalised by Satyanarayanan
et al.~\cite{satyanarayanan2009case}, who characterised a cloudlet
as a resource-rich node placed one network hop away from mobile
users. By co-locating cloudlets with Wi-Fi access points or campus
gateways, latency-sensitive tasks such as augmented reality and
real-time video analysis can be processed in milliseconds rather
than the hundreds of milliseconds typical of cloud round-trips.
Cloudlets offer a practical middle ground between the resource
abundance of the cloud and the proximity of on-device processing.

\subsection{Fog Computing}

Bonomi et al.~\cite{bonomi2012fog} introduced fog computing as a
model that extends cloud services to the network edge. Fog
nodes---implemented in routers, switches, and IoT gateways---aggregate,
filter, and locally analyse data streams before forwarding selected
results upstream. A layered arrangement lets fog gateways handle
lightweight preprocessing while cloudlets manage more demanding
workloads, substantially reducing WAN bandwidth consumption.

\subsection{Multi-Access Edge Computing}

The European Telecommunications Standards Institute (ETSI)
standardised Multi-Access Edge Computing (MEC), placing compute
resources at cellular and network access points~\cite{etsi2016mec}.
MEC supports 5G applications with strict latency requirements---
vehicular control, online gaming, and industrial automation---without
incurring WAN round-trip delays. The ecosystem also exposes
location-aware APIs so applications can be scheduled in proximity
to the serving base station.

% ─────────────────────────────────────────────────────────
\section{Load Balancing Fundamentals and Strategies}
\label{sec:lit_lb}

Load balancing distributes computing tasks across available
resources to improve throughput, latency, and utilisation. In
edge environments, effective load management is especially
important because cloudlets have limited capacity and demand
patterns are highly variable.

\subsection{Static Load Balancing}

Static algorithms apply predetermined assignment rules without
regard to current system state.

\begin{itemize}
  \item \textbf{Round Robin (RR):} Tasks are assigned to nodes
  cyclically, irrespective of node capacity or current load. RR
  is simple and imposes negligible overhead but ignores workload
  heterogeneity.

  \item \textbf{Weighted Round Robin (WRR):} Each node receives a
  capacity-proportional weight, so more powerful nodes handle more
  requests. WRR improves on RR for heterogeneous settings but
  cannot adapt when runtime capacity changes.

  \item \textbf{Hash-Based Methods:} Tasks are assigned by hashing
  a request attribute such as a client identifier, achieving
  session affinity and stateless distribution. Skewed key
  distributions can, however, produce significant imbalance.
\end{itemize}

Static methods are predictable and easy to implement but lack the
flexibility to react to changing workload patterns or node
failures, making them unsuitable for dynamic edge environments.

\subsection{Dynamic Load Balancing}

Dynamic policies observe the system state continuously and adjust
assignment decisions in real time.

\begin{itemize}
  \item \textbf{Least Connection (LC):} New requests go to the
  server with the fewest active connections. LC suits long-lived
  sessions but may be suboptimal for short-lived requests where
  connection overhead dominates.

  \item \textbf{Response-Time Based:} The scheduler tracks actual
  server delay and directs requests to the fastest node. This
  delivers strong per-request performance but is sensitive to
  rapid latency fluctuations and requires continuous measurement.

  \item \textbf{Resource-Based Balancing (LBRO):} The LBRO
  framework~\cite{nayyer2022lbro} computes a Composite Resource
  Index (CRI) for each cloudlet as a weighted sum of available
  CPU, RAM, and queue capacity, routing tasks to the highest-CRI
  node. LBRO outperforms round-robin on drop rate and latency
  under moderate load, but is purely reactive---it uses only the
  current observed state and has no mechanism to anticipate
  congestion or learn from past decisions.
\end{itemize}

Dynamic techniques respond faster than static ones but remain
reactive and require monitoring infrastructure. Threshold-based
offloading policies, which trigger cloud forwarding when
utilisation crosses a fixed limit, share the same reactiveness
and require per-deployment manual tuning~\cite{mach2017mobile}.

% ─────────────────────────────────────────────────────────
\section{Machine Learning-Based Load Balancing}
\label{sec:lit_ml}

\subsection{Supervised Classification Approaches}

Supervised learning methods have been used to incorporate task
characteristics and historical patterns into scheduling decisions.
Several studies trained Random Forest (RF) or gradient-boosted
classifiers on labelled traces to predict offloading destinations,
using features such as task size, estimated execution time, and
current cloudlet load~\cite{wang2019task}. These classifiers
achieve reasonable accuracy when training and deployment
distributions match, but depend on labelled data and cannot adapt
online as workload statistics change.

\subsection{Predictive Pre-Allocation}

The Pred-LB framework~\cite{predlb2025} forecasts user activity
patterns and pre-allocates resources accordingly, reducing reactive
overloading during demand surges. While effective for periodic or
predictable demand, such methods are less reliable when task
arrivals are driven by irregular events that are difficult to
anticipate from activity logs.

A shared limitation of supervised approaches is that the
scheduling objective is defined implicitly through the labelling
process rather than directly optimised. If the training labels
were produced by a suboptimal policy, the learned classifier
inherits those shortcomings. Reinforcement learning avoids this
by directly optimising cumulative reward.

% ─────────────────────────────────────────────────────────
\section{Deep Reinforcement Learning for MEC Scheduling}
\label{sec:lit_drl}

\subsection{Value-Based Methods}

The foundational DeepRM work of Mao et al.~\cite{mao2016resource}
demonstrated that a policy gradient agent can learn job scheduling
policies that outperform conventional heuristics on makespan
metrics. Liu et al.~\cite{liu2022edgelb} proposed DRL-EdgeLB, a
DQN-based scheme where the agent selects a task destination based
on per-server queue lengths and CPU utilisation. Experiments showed
improvements over CRI-based heuristics under high load. However,
DRL-EdgeLB's state representation is purely reactive and includes
no forecast of future load, limiting its ability to act
proactively when a cloudlet is trending toward congestion.

\subsection{Multi-Agent and Policy Gradient Methods}

Zhao et al.~\cite{zhao2022madrl} explored MADRL-MEC, a multi-agent
formulation in which each edge server runs its own agent
coordinated through a shared critic. While the multi-agent approach
can tailor policies to individual server characteristics, it
introduces additional complexity in training stability and credit
assignment that may be unnecessary when a centralised broker can
observe the full system state.

PPO~\cite{schulman2017ppo} has also been applied to joint
offloading and resource allocation in MEC
settings~\cite{chen2019optimized}. As an on-policy method, PPO
must discard collected experience after each policy update, making
it less sample-efficient than off-policy algorithms such as DDQN
in simulation-based settings where environment interaction is
inexpensive.

% ─────────────────────────────────────────────────────────
\section{Predictive Scheduling with Recurrent Neural Models}
\label{sec:lit_lstm}

\subsection{LSTM and GRU for Load Forecasting}

LSTM networks~\cite{hochreiter1997lstm} address the vanishing
gradient problem of plain recurrent networks and have been widely
used to model server CPU and memory utilisation. Qiu et
al.~\cite{qiu2020lstm} showed that LSTM-based predictors track
non-stationary utilisation patterns with lower mean absolute error
than autoregressive baselines, particularly when utilisation
exhibits diurnal periodicities. The GRU~\cite{cho2014gru}, a
lighter alternative using reset and update gates, achieves
comparable accuracy on shorter sequences. Hybrid GRU-LSTM models,
where a GRU layer extracts local temporal features and an LSTM
layer captures longer dependencies, have shown advantages on
server load traces where both short-term fluctuations and
longer-horizon trends are present.

\subsection{Combining Forecasting with Reinforcement Learning}

Zhang et al.~\cite{zhang2021joint} proposed a joint prediction and
offloading framework in which an LSTM forecasts channel quality and
task arrival rates, feeding these predictions to a DQN agent.
The approach demonstrated that augmenting the RL state with
predicted quantities improves scheduling quality in dynamic
wireless environments. Crucially, the agent's value function can
exploit predictions only when they are part of the state it is
trained on from the beginning---not when appended as a post-hoc
feature at evaluation time.

% ─────────────────────────────────────────────────────────
\section{Comparison of Related Work}
\label{sec:lit_comparison}

Table~\ref{tab:related_summary} summarises key attributes of
the approaches discussed and positions DRL-LSTM-LBRO relative
to existing work.

\begin{table}[htbp]
\centering
\caption{Comparison of Edge Load Balancing and Scheduling Approaches}
\label{tab:related_summary}
\footnotesize
\renewcommand{\arraystretch}{1.15}
\begin{tabular}{|p{2.8cm}|p{2.5cm}|p{1.1cm}|p{1.5cm}|p{1.3cm}|p{1.3cm}|}
\hline
\textbf{Reference} &
\textbf{Technique} &
\textbf{Adaptive} &
\textbf{Multi-Resource} &
\textbf{Proactive} &
\textbf{Scalable} \\
\hline
Static RR / WRR & Fixed rules & No & Limited & No & High \\
\hline
LC / Response-Time & Real-time metrics & Yes & Partial & No & Medium \\
\hline
LBRO \cite{nayyer2022lbro} & Weighted CRI & Yes & Yes & No & Medium \\
\hline
DeepRM \cite{mao2016resource} & Policy gradient & Yes & Yes & Partial & Low \\
\hline
DRL-EdgeLB \cite{liu2022edgelb} & DQN, reactive & Yes & Yes & No & Medium \\
\hline
MADRL-MEC \cite{zhao2022madrl} & Multi-agent DRL & Yes & Yes & No & Low \\
\hline
Pred-LB \cite{predlb2025} & Activity prediction & Yes & Partial & Yes & Medium \\
\hline
PPO \cite{schulman2017ppo} & Policy gradient & Yes & Yes & Partial & Medium \\
\hline
Zhang et al.\ \cite{zhang2021joint} & LSTM + DQN & Yes & Partial & Yes & Low \\
\hline
\textbf{DRL-LSTM-LBRO (ours)} & \textbf{GRU-LSTM + DDQN} & \textbf{Yes} & \textbf{Yes} & \textbf{Yes} & \textbf{High} \\
\hline
\end{tabular}
\end{table}

% ─────────────────────────────────────────────────────────
\section{Research Gaps and Motivation}
\label{sec:lit_gaps}

The literature survey reveals four key gaps that the present work
addresses.

\begin{enumerate}

  \item \textbf{Reactive state representations in DRL schedulers.}
  Existing schedulers such as DRL-EdgeLB and MADRL-MEC use only
  the current observed resource state. A cloudlet that appears
  lightly loaded at decision time may already be trending toward
  saturation. Incorporating near-future load predictions into the
  RL state would allow the agent to act preventively.

  \item \textbf{Isolated treatment of forecasting and scheduling.}
  Where prediction and scheduling have been combined, forecasts are
  produced by a separately trained model and passed to a fixed
  heuristic. The scheduling policy is not trained with predictions
  in its state, creating a distribution mismatch that can degrade
  performance. An end-to-end approach that includes LSTM
  predictions in the agent's state from the start is preferable.

  \item \textbf{Single-objective or incomplete reward functions.}
  Many DRL schedulers optimise a single metric or omit practically
  important criteria such as load imbalance or overload penalties.
  Incomplete rewards are susceptible to policy collapse, where the
  agent routes all traffic to the most powerful node, degrading
  fairness and accelerating saturation under variable load.

  \item \textbf{Limited scalability analysis.}
  Most published DRL schedulers are evaluated on a fixed small
  number of cloudlets (typically two to four). Whether the learned
  policy generalises as the cloudlet count grows is rarely tested,
  despite being important for real MEC deployments.

\end{enumerate}

The DRL-LSTM-LBRO framework proposed in this thesis closes all
four gaps. Per-cloudlet GRU-LSTM predictors are incorporated into
the DDQN agent's state from the start of training. The reward
function covers six criteria including explicit imbalance and
overload penalties. The framework is evaluated on both
three-cloudlet and six-cloudlet topologies to verify scalability.
